\chapter{参考文献}
\begin{enumerate}
\item 周建伟：《微分几何讲义》，科学出版社。主参考。出版社试读的版本信息为2010年第一版，ISBN 978-7-03-028107-4；本地DjVu的版次与页码仍须解码核对。
\item S. S. Chern, W. H. Chen, K. S. Lam: \textit{Lectures on Differential Geometry}。本地文件名虽为中文，但封面及内文为英文；不应误记为中文教材。
\item 梁灿彬、周彬：《微分几何入门与广义相对论》。本地扫描件首页标注上册、第二版；是否还包含下册需继续核对，不以文件名推断覆盖范围。
\item Marco Gualtieri: \textit{Geometry and Topology, 1300Y}, 2008--2009。补充参考：4.1的形式与同伦，4.3--4.4的Mayer--Vietoris及球面计算（印刷/PDF页42--44）；\href{https://www.math.toronto.edu/mgualt/MAT1300/1300\%20Lecture\%20notes.pdf}{多伦多大学公开讲义}。本书重新组织证明并固定自己的差映射符号。
\item Eckhard Meinrenken: \textit{Riemannian Geometry}, University of Toronto, Spring 2002。补充参考12--13节的度量与存在性、15节的联络与基本定理、16--17节的指数映射/完备性及18--19节的曲率与向量丛；\href{https://www.math.toronto.edu/mein/teaching/LectureNotes/rieall.pdf}{公开讲义}。本书以曲率为零定义一般联络的平坦性，并单列无挠时的平行坐标结论。
\item Clara L\"oh: \textit{Differential Geometry I}, Universit\"at Regensburg, Wintersemester 2020/21。补充核对3.5节第二变分、Jacobi场、共轭与内点共轭后的非最短性（印刷159--169/PDF167--177）；\href{https://loeh.app.uni-regensburg.de/teaching/diffgeo_ws2021/lecture_notes.pdf}{官方讲义}。本书另补指标引理的矩阵证明及负指标场的光滑化细节。
\item Ben Andrews: \textit{Lectures on Differential Geometry}, Australian National University。补充核对第17--19讲（印刷165--180）的法丛、相容方程及活动标架；\href{https://maths-people.anu.edu.au/~andrews/DG/}{官方课程与分讲PDF}。本书另在一般黎曼环境推导曲率项，并以自身符号重写，局部浸入重建和体积变分由前章工具补证。
\end{enumerate}
